\section{Introduction}
\label{sec:intro}

According to the World Health Organization, falls are the second
leading cause of accidental injury deaths globally, with the vast
majority occurring in indoor domestic environments~\cite{who2021falls}.
The scale of this problem is significant: the United States Centers
for Disease Control and Prevention estimates that falls cause over
36,000 deaths and 8~million emergency department visits annually in
the US alone. Beyond falls, lacerations from sharp objects--knives,
scissors, and broken glass--represent a further substantial source
of domestic injury, particularly among children under five and adults
over sixty-five. Furniture and clutter obstacles contribute additional
collision and entrapment injuries across home, hospital, and retail
environments.

The development of intelligent domestic systems creates an urgent
need for vision systems capable of parsing indoor scenes in terms
of their \textit{safety consequences} rather than purely their
semantic content. Assistive robots navigating hospital corridors,
smart monitoring systems for elderly residents, autonomous
wheelchairs operating in homes, and occupational safety inspection
systems all require pixel-level knowledge of which regions are safe
to traverse and which present hazards. A system that cannot
distinguish a wet slippery floor from a dry safe floor, or that
classifies a knife on the ground identically to a harmless book,
is fundamentally inadequate for safety-critical deployment.

Semantic segmentation--the task of assigning a class label to every
pixel in an image--provides the natural representational framework
for this problem. However, applying existing segmentation architectures
to indoor hazard understanding exposes two fundamental gaps that
current methods do not address.

\noindent\textbf{Gap~1: Missing hazard-specific annotations.}
No existing segmentation benchmark defines pixel classes aligned
with indoor injury risk. ADE20K~\cite{zhou2017scene}, the largest
indoor scene parsing benchmark, provides 150~semantic categories
but treats all floor types identically and has no concept of
obstacle severity or laceration risk. The COCO
dataset~\cite{lin2014coco} provides instance segmentation for
everyday objects including knives and scissors, but does not frame
these in terms of safety consequence. NYU Depth V2~\cite{silberman2012indoor}
and SUN~RGB-D~\cite{song2015sun} similarly provide generic indoor
annotations without safety-critical class definitions.

\noindent\textbf{Gap~2: Single-level simultaneous prediction.}
Every existing semantic segmentation architecture--U-Net~\cite{ronneberger2015unet},
DeepLab~\cite{chen2017deeplab}, PSPNet~\cite{zhao2017pyramid},
SegFormer~\cite{xie2021segformer}--predicts all semantic classes
simultaneously from a single decoder pass over encoder features.
This design provides no mechanism to first answer whether a region
is dangerous before classifying its specific hazard type, missing
a natural two-level hierarchy inherent to safety reasoning. A
human inspector, when assessing indoor hazard risk, naturally first
identifies which regions deserve attention (coarse binary detection)
and then classifies the specific hazard type (fine-grained recognition).
No existing architecture models this reasoning structure.

This paper addresses both gaps through a unified architectural
and data contribution. For Gap~1, we introduce the first
hazard-specific segmentation taxonomy and contribute a manually
annotated indoor hazard dataset. For Gap~2, we propose DSPHN,
which integrates two previously separate architectural paradigms
into a single model:

\begin{enumerate}
\item \textbf{DAHD dual-stream decoding:} Parallel Appearance and
Context decoders with severity-guided Fusion Gates disentangle the
two key aspects of hazard recognition: local visual appearance cues
(\textit{what} does a hazard look like?) and spatial contextual cues
(\textit{where} does a hazard occur relative to its surroundings?).

\item \textbf{PHRN hierarchical inference:} A two-stage conditional
design where Stage~1 first produces a binary safe/unsafe prediction
from the fused dual-stream features, and Stage~2 receives this
binary map as a structural input to perform fine-grained four-class
hazard classification.
\end{enumerate}

The specific contributions of this paper are:

\begin{itemize}
\item \textbf{DSPHN: Novel Integrated Architecture.} The first
segmentation architecture to combine dual-stream appearance-context
disentanglement with hierarchical safety-aware conditional inference
for indoor hazard segmentation from a single RGB image. The
combination produces emergent benefits: Stage~1 maps derived from
dual-stream features are more accurate than single-decoder Stage~1
maps, providing better conditioning for Stage~2 and explaining why
DSPHN outperforms both parent architectures.

\item \textbf{Controlled Ablation Validation.} All three variants
(DAHD-only, PHRN-only, DSPHN) are trained and evaluated under
identical conditions. DSPHN achieves mIoU~0.4198 versus
0.3894~(DAHD-only) and 0.3793~(PHRN-only), demonstrating that
integration is quantitatively superior to either component alone.

\item \textbf{Indoor Hazard Taxonomy.} A four-class safety-oriented
taxonomy--Background, Safe~Floor, Obstacle, Sharp~Objects--defined
by real-world injury consequence rather than semantic identity,
enabling training signal directly aligned with safety outcomes.

\item \textbf{Manual Hazard Annotation Dataset.} 100~ADE20K indoor
images manually annotated at pixel level using Make~Sense~AI with
three hazard categories, constituting the first annotation set
combining floor safety, obstacle, and sharp object categories for
indoor segmentation.

\item \textbf{COCO Sharp Object Augmentation.} Integration of
186~COCO knife and scissors instances without floor-level
spatial constraint, motivated by the observation that sharp
objects constitute laceration hazards regardless of their
vertical position in a scene.
\end{itemize}

The remainder of this paper is organised as follows.
Section~\ref{sec:related} reviews related work in semantic
segmentation, multi-stage architectures, attention mechanisms,
loss functions, and indoor safety understanding.
Section~\ref{sec:method} describes the DSPHN architecture in
full mathematical detail.
Section~\ref{sec:data} presents the dataset construction and
annotation methodology.
Section~\ref{sec:experiments} reports experimental results
including ablation study, per-class analysis, training dynamics,
qualitative evaluation, and extended discussion.
Section~\ref{sec:conclusion} concludes with future directions.

\section{Related Work}
\label{sec:related}

\subsection{Semantic Segmentation Architectures}

The foundational work of Long~\etal~\cite{long2015fully} established
fully convolutional networks (FCNs) for dense pixel prediction,
replacing classification heads with upsampling paths to enable
end-to-end training for per-pixel label assignment. The
encoder-decoder paradigm was significantly refined by
Ronneberger~\etal~\cite{ronneberger2015unet}, whose U-Net
architecture introduced symmetric skip connections between encoder
and decoder stages to preserve high-frequency spatial details
lost during downsampling, originally demonstrated on biomedical
image segmentation. U-Net's elegant design has since become
a backbone architecture across diverse domains including remote
sensing, medical imaging, autonomous driving, and indoor scene
understanding.

DeepLab~\cite{chen2017deeplab} addressed the need for multi-scale
context through atrous (dilated) convolution and Atrous Spatial
Pyramid Pooling (ASPP), enabling the network to capture information
at multiple receptive field scales without increasing the number
of parameters. PSPNet~\cite{zhao2017pyramid} employed a Pyramid
Pooling Module to aggregate global scene context through
multi-resolution pooling. Subsequent transformer-based architectures
have achieved state-of-the-art results on benchmark datasets:
SegFormer~\cite{xie2021segformer} combines a hierarchical Mix
Transformer encoder with a lightweight all-MLP decoder, while
Swin-UNet~\cite{cao2022swinunet} adapts the Swin Transformer
for pure-transformer medical image segmentation.

A critical observation underlying DSPHN is that \textit{all} of
these architectures share a common design principle: a single decoder
that predicts all semantic classes simultaneously from encoded
features. This design treats hazard and non-hazard regions
identically, providing no hierarchical structure for safety-aware
reasoning. DSPHN departs from this paradigm by introducing a
two-stage conditional decoder.

\subsection{Multi-Stream and Dual-Decoder Architectures}

Dual-decoder and multi-stream designs have been primarily explored
in RGB-D segmentation, where separate processing branches handle
the RGB and depth input modalities~\cite{chen2020bi}. The
bi-directional modality fusion approach of Chen~\etal~\cite{chen2020bi}
uses cross-modal attention to propagate features between the two
streams. In the single-modality RGB setting, dilated convolutions
have been used in DeepLab variants to capture multi-scale context,
but within a single decoder pathway.

DSPHN applies parallel Appearance and Context decoders to a
\textit{single RGB image}. The dual-stream design is not driven
by two input modalities but by the semantic distinction between
two aspects of hazard recognition that benefit from different
inductive biases: standard convolutions for local texture pattern
learning (Appearance Decoder) and dilated convolutions for
spatial relationship modelling over larger receptive fields
(Context Decoder). This is a qualitatively different motivation
from RGB-D dual-stream designs.

\subsection{Coarse-to-Fine and Multi-Stage Architectures}

Coarse-to-fine prediction strategies have been applied across
multiple vision tasks. Eigen~\etal~\cite{eigen2014depth} used
a two-stage coarse-to-fine network for depth estimation, where a
global prediction is refined at higher resolution. Cascade
R-CNN~\cite{cai2018cascade} employs sequential bounding box
regression stages for object detection, where each stage receives
refined proposals from the previous. In semantic segmentation,
label refinement networks~\cite{islam2017label} refine
coarse predictions at progressively finer resolutions.

DSPHN's two-stage design is qualitatively different from all of
these. Standard coarse-to-fine approaches refine the \textit{same
output type} (depth, bounding box, segmentation map) at increasing
resolution. In DSPHN, Stage~1 and Stage~2 answer
\textit{categorically different questions}: Stage~1 performs binary
safe/unsafe detection, while Stage~2 performs four-class hazard
classification. Stage~1 output is a different type of prediction
entirely, not a coarser version of Stage~2 output.

\subsection{Attention Mechanisms in Segmentation}

Attention mechanisms have been widely adopted to improve feature
selectivity. The Squeeze-and-Excitation Network~\cite{hu2018squeeze}
introduced channel-wise attention via global average pooling to
recalibrate channel responses. Non-local networks~\cite{wang2018nonlocal}
modelled long-range spatial dependencies through self-attention.
Attention U-Net~\cite{oktay2018attention} proposed soft spatial
attention gates on skip connections, using decoder gating signals
to suppress irrelevant encoder activations.

DSPHN's Severity Fusion Gates differ from these in three ways:
(1) they blend two \textit{complete decoder pathways} rather than
modulating individual skip connections; (2) they are explicitly
initialised from domain hazard severity scores providing a
principled starting point for fusion learning; and (3) they learn
when to trust appearance versus spatial context for each pixel,
rather than learning which encoder features to suppress.

\subsection{Loss Functions for Class Imbalance}

Class imbalance is pervasive in semantic segmentation.
Class-weighted Cross-Entropy assigns higher penalties to
underrepresented classes but optimises pixel-level accuracy.
Dice Loss~\cite{milletari2016vnet} directly optimises the F1
overlap coefficient, exhibiting less sensitivity to class frequency
imbalance. Focal Loss~\cite{lin2017focal} dynamically downweights
well-classified easy examples to focus training on hard, typically
minority, pixels. Taghanaki~\etal~\cite{taghanaki2019combo}
demonstrated that combining Cross-Entropy and Dice Loss consistently
outperforms either alone across imbalanced segmentation tasks.
DSPHN extends this combination with an additional binary BCE term
providing direct Stage~1 supervision, ensuring the binary hazard
map is accurate before Stage~2 depends on it.

\subsection{Indoor Scene Understanding and Safety}

General indoor segmentation benchmarks (ADE20K~\cite{zhou2017scene},
NYU Depth V2~\cite{silberman2012indoor}, SUN RGB-D~\cite{song2015sun})
provide rich semantic annotations for generic indoor parsing but
define no hazard-specific classes and encode no safety consequence
information. Prior work addressing specific narrow indoor hazards
includes staircase detection~\cite{tummala2021stair} using depth
sensors, IoT-based wet floor alert systems, and object detection
for robotic obstacle avoidance. These approaches either require
specialised sensors unavailable in many deployments, or address
only a single hazard type in isolation.

To our knowledge, no prior work formulates general indoor hazard
understanding as a unified semantic segmentation problem with
hazard-class-aware supervision over floor safety, obstacle risk,
and sharp object detection simultaneously from a standard RGB
camera. This represents the application gap that DSPHN addresses,
combined with the architectural novelty of the dual-stream
hierarchical design.
